% Part: incompleteness
% Chapter: introduction
% Section: definitions

\documentclass[../../../include/open-logic-section]{subfiles}

\begin{document}

\olfileid[id]{inc}{int}{def}

\olsection{Definisi}

Untuk melaksanakan proyek Hilbert dalam memformalkan matematika serta
menunjukkan bahwa formalisasi tersebut konsisten dan lengkap, hal pertama yang
harus dilakukan ialah memilih bahasa, kerangka logis, dan sistem aksioma.
Untuk keperluan kita, andaikan matematika dapat diformalkan dalam bahasa orde
pertama, yakni terdapat suatu himpunan !!{constant}s, !!{function}s, dan
!!{predicate}s yang, bersama penghubung dan kuantor logika orde pertama,
memungkinkan kita menyatakan klaim-klaim matematika. Kebanyakan orang sepakat
bahwa bahasa semacam itu ada: bahasa teori himpunan, dengan $\in$ sebagai
satu-satunya simbol nonlogis. Sekilas, klaim bahwa bahasa sesederhana itu
begitu ekspresif tentu tampak sangat tidak masuk akal, dan diperlukan banyak
kerja untuk menetapkan bahwa hampir seluruh matematika dapat dinyatakan dengan
kosakata yang sangat sederhana ini. Agar pembahasan tetap sederhana, untuk
sementara kita batasi pada aritmetika, yakni bagian matematika yang hanya
berurusan dengan bilangan asli~$\Nat$. Bahasa yang alami untuk menyatakan
fakta aritmetika ialah~$\Lang L_A$. $\Lang L_A$ memuat satu !!{predicate}
dua-tempat~$<$, satu !!{constant}~$\Obj 0$, satu !!{function}
satu-tempat~$\prime$, serta dua !!{function}s dua-tempat~$+$ dan~$\times$.

\begin{defn}
Suatu himpunan !!{sentence}s~$\Gamma$ adalah sebuah \emph{teori} jika tertutup
terhadap perikutan, yakni jika $\Gamma = \Setabs{!A}{\Gamma \Entails !A}$.
\end{defn}

Ada dua cara mudah untuk menentukan teori. Cara pertama ialah menjadikannya
sebagai himpunan !!{sentence}s yang benar dalam suatu !!{structure}. Sebagai
contoh, perhatikan !!{structure} bagi $\Lang L_A$ yang !!{domain}-nya
ialah~$\Nat$ dan semua simbol nonlogisnya ditafsirkan sebagaimana mestinya.

\begin{defn}\ollabel{def:standard-model}
\emph{Model standar aritmetika} ialah !!{structure}~$\Struct N$ yang
didefinisikan sebagai berikut:
\begin{enumerate}
\item $\Domain N = \Nat$
\item $\Assign{\Obj 0}{N} = 0$
\item $\Assign{\Obj \prime}{N}(n) = n + 1$ untuk semua $n \in \Nat$
\item $\Assign{\Obj +}{N}(n, m) = n + m$ untuk semua $n, m \in \Nat$
\item $\Assign{\Obj \times}{N}(n, m) = n\cdot m$ untuk semua $n, m \in \Nat$
\item $\Assign{\Obj <}{N} = \Setabs{\tuple{n, m}}{n \in \Nat, m \in
  \Nat, n < m}$
\end{enumerate}
\end{defn}

Perhatikan perbedaan antara `$\times$' dan `$\cdot$': $\times$ adalah simbol
dalam bahasa aritmetika. Tentu saja, simbol itu dipilih agar mengingatkan kita
pada perkalian, tetapi $\times$ bukan operasi perkalian, melainkan simbol
fungsi dua-tempat (secara resmi, $\Obj f^2_1$). Sebaliknya,
$\cdot$~\emph{adalah} fungsi perkalian biasa. Ketika melihat sesuatu seperti
$n \cdot m$, yang dimaksud adalah hasil kali bilangan $n$ dan~$m$; ketika
melihat sesuatu seperti $x \times y$, kita sedang membicarakan suatu suku
dalam bahasa aritmetika. Dalam model standar, simbol fungsi~$\times$
ditafsirkan sebagai fungsi $\cdot$ pada bilangan asli. Untuk penjumlahan, kita
menggunakan $+$ sebagai simbol fungsi bahasa aritmetika sekaligus fungsi
penjumlahan pada bilangan asli. Di sini, konteks menentukan mana yang dimaksud.

\begin{defn}
Teori \emph{aritmetika sejati} ialah himpunan !!{sentence}s yang dipenuhi dalam
model standar aritmetika, yakni
\[
\Th{TA} = \Setabs{!A}{\Sat{N}{!A}}.
\]
\end{defn}

$\Th{TA}$ adalah suatu teori, sebab setiap kali $\Th{TA} \Entails !A$,
$!A$ dipenuhi dalam setiap !!{structure} yang memenuhi~$\Th{TA}$. Karena
$\Sat{N}{\Th{TA}}$, diperoleh $\Sat{N}{!A}$, sehingga $!A \in \Th{TA}$.

Cara lain untuk menentukan teori~$\Gamma$ ialah menjadikannya sebagai
himpunan !!{sentence}s yang merupakan konsekuensi logis dari suatu himpunan
kalimat~$\Gamma_0$.
Dalam hal ini, $\Gamma$ merupakan ``penutupan'' $\Gamma_0$ terhadap
perikutan. Menentukan teori dengan cara ini hanya menarik jika $\Gamma_0$
ditentukan secara eksplisit, misalnya jika !!{element}s~$\Gamma_0$ didaftar.
Sekurang-kurangnya, $\Gamma_0$ harus dapat diputuskan, yakni harus ada uji
yang dapat dihitung untuk menentukan kapan !!a{sentence} merupakan anggota
$\Gamma_0$ dan kapan bukan. Kita menyebut !!{sentence}s dalam~$\Gamma_0$
sebagai \emph{aksioma} bagi~$\Gamma$, dan mengatakan bahwa $\Gamma$
\emph{diaksiomatisasi} oleh~$\Gamma_0$.

\begin{defn}
Suatu teori~$\Gamma$ \emph{diaksiomatisasi} oleh~$\Gamma_0$ jika dan hanya
jika
\[
\Gamma = \Setabs{!A}{\Gamma_0 \Entails !A}
\]
\end{defn}

\begin{defn}
Teori $\Th{Q}$ yang diaksiomatisasi oleh kalimat-kalimat berikut dikenal
sebagai ``$\Th{Q}$ Robinson'' dan merupakan teori aritmetika yang sangat
sederhana.
\begin{align*}
& \lforall[x][\lforall[y][(\eq[x'][y'] \lif \eq[x][y])]] \tag{$!Q_1$}\\
& \lforall[x][\eq/[\Obj 0][x']] \tag{$!Q_2$}\\
& \lforall[x][(\eq[x][\Obj 0] \lor \lexists[y][\eq[x][y']])] \tag{$!Q_3$}\\
& \lforall[x][\eq[(x + \Obj 0)][x]] \tag{$!Q_4$}\\
& \lforall[x][\lforall[y][\eq[(x + y')][(x + y)']]] \tag{$!Q_5$}\\
& \lforall[x][\eq[(x \times \Obj 0)][\Obj 0]] \tag{$!Q_6$}\\
& \lforall[x][\lforall[y][\eq[(x \times y')][((x \times y) + x)]]] \tag{$!Q_7$}\\
& \lforall[x][\lforall[y][(x < y \liff \lexists[z][\eq[(z' + x)][y]])]] \tag{$!Q_8$}
\end{align*}
Himpunan !!{sentence}s $\{!Q_1, \dots, !Q_8\}$ adalah aksioma-aksioma
$\Th{Q}$, sehingga $\Th{Q}$ terdiri atas semua !!{sentence}s yang merupakan
konsekuensi logis dari aksioma-aksioma tersebut:
\[
\Th{Q} = \Setabs{!A}{\{!Q_1, \dots, !Q_8\} \Entails !A}.
\]
\end{defn}

\begin{defn}
Andaikan $!A(x)$ adalah !!a{formula} dalam $\Lang L_A$ dengan variabel bebas
$x$ dan $y_1$, \dots, $y_n$. Maka setiap !!{sentence} berbentuk
\[
\lforall[y_1][\dots\lforall[y_n][((!A(\Obj 0) \land \lforall[x][(!A(x)
\lif !A(x'))]) \lif \lforall[x][!A(x)])]]
\]
merupakan suatu contoh \emph{skema induksi}.

\emph{Aritmetika Peano}~$\Th{PA}$ ialah teori yang diaksiomatisasi oleh
aksioma-aksioma $\Th{Q}$ bersama semua contoh skema induksi.
\end{defn}

\begin{explain}
Setiap contoh skema induksi benar dalam~$\Struct{N}$. Hal ini paling mudah
dilihat jika !!{formula}~$!A$ hanya memiliki satu !!{variable} bebas~$x$.
Maka $!A(x)$ mendefinisikan himpunan bagian~$X_{!A}$ dari~$\Nat$
dalam~$\Struct{N}$. $X_{!A}$ adalah himpunan semua~$n \in \Nat$ sedemikian
sehingga $\Sat{N}{!A(x)}[s]$ ketika $s(x) = n$. Contoh skema induksi yang
bersesuaian ialah
\[
((!A(\Obj 0) \land \lforall[x][(!A(x) \lif !A(x'))]) \lif
  \lforall[x][!A(x)]).
\]
Jika antesedennya benar dalam~$\Struct{N}$, maka $0 \in X_{!A}$ dan, setiap
kali $n \in X_{!A}$, demikian pula $n+1$. Karena $0 \in X_{!A}$,
kita memperoleh $1 \in X_{!A}$. Dari $1 \in X_{!A}$ kita memperoleh
$2 \in X_{!A}$, dan seterusnya. Jadi untuk setiap $n \in \Nat$,
$n \in X_{!A}$. Namun, ini berarti bahwa $\lforall[x][!A(x)]$ dipenuhi
dalam~$\Struct{N}$.
\end{explain}

Baik $\Th{Q}$ maupun $\Th{PA}$ merupakan teori yang diaksiomatisasi.
Pertanyaan besarnya ialah: seberapa kuat kedua teori itu? Misalnya, dapatkah
$\Th{PA}$ membuktikan semua kebenaran tentang~$\Nat$ yang dapat dinyatakan
dalam~$\Lang L_A$? Secara khusus, apakah aksioma-aksioma $\Th{PA}$
menyelesaikan semua pertanyaan yang dapat dirumuskan dalam~$\Lang L_A$?

Cara lain untuk menyatakannya ialah bertanya: apakah $\Th{PA} = \Th{TA}$?
$\Th{TA}$ jelas membuktikan (yakni memuat) semua kebenaran tentang~$\Nat$,
dan menyelesaikan semua pertanyaan yang dapat dirumuskan dalam~$\Lang L_A$,
sebab jika $!A$ adalah !!a{sentence} dalam $\Lang L_A$, maka
$\Sat{N}{!A}$ atau $\Sat{N}{\lnot !A}$; jadi $\Th{TA} \Entails !A$ atau
$\Th{TA} \Entails \lnot !A$. Teori semacam itu disebut
\emph{!!{complete}}.

\begin{defn}
Suatu teori $\Gamma$ \emph{!!{complete}} jika dan hanya jika, untuk setiap
!!{sentence}~$!A$ dalam bahasanya, $\Gamma \Entails !A$ atau $\Gamma
\Entails \lnot !A$.
\end{defn}

\begin{explain}
Berdasarkan Teorema Kelengkapan, $\Gamma \Entails !A$ jika dan hanya jika
$\Gamma \Proves !A$; jadi $\Gamma$ lengkap jika dan hanya jika, untuk setiap
!!{sentence}~$!A$ dalam bahasanya, $\Gamma \Proves !A$ atau $\Gamma
\Proves \lnot !A$.
\end{explain}

Pertanyaan lain yang muncul ialah: adakah prosedur komputasional untuk
menguji apakah !!a{sentence} terdapat dalam~$\Th{TA}$, $\Th{PA}$, atau
bahkan hanya dalam~$\Th{Q}$? Kita dapat membuat pertanyaan ini lebih tepat
dengan mendefinisikan kapan suatu himpunan (misalnya himpunan !!{sentence}s)
bersifat !!{decidable}.

\begin{defn}
Suatu himpunan $X$ \emph{!!{decidable}} jika dan hanya jika terdapat prosedur
komputasional yang, dengan masukan~$x$, mengembalikan~$1$ jika $x \in X$ dan
$0$ jika tidak.
\end{defn}

Jadi pertanyaan kita menjadi: apakah $\Th{TA}$ ($\Th{PA}$, $\Th{Q}$)
!!{decidable}?

Jawaban atas semua pertanyaan ini ialah: tidak. Tidak satu pun teori tersebut
dapat diputuskan. Namun, fenomena ini tidak khusus bagi teori-teori tersebut.
Sesungguhnya, \emph{setiap} teori yang memenuhi syarat tertentu terkena hasil
yang sama. Salah satu syarat yang dipenuhi $\Th{Q}$ dan $\Th{PA}$ ialah
bahwa keduanya diaksiomatisasi oleh !!a{decidable} himpunan aksioma.

\begin{defn}
Suatu teori \emph{!!{axiomatizable}} jika diaksiomatisasi oleh
!!a{decidable} himpunan aksioma.
\end{defn}

\begin{ex}
Setiap teori yang diaksiomatisasi oleh himpunan !!{sentence}s berhingga adalah
!!{axiomatizable}, sebab setiap himpunan berhingga !!{decidable}. Jadi,
misalnya, $\Th{Q}$ !!{axiomatizable}.

Teori yang diaksiomatisasi secara skematis seperti~$\Th{PA}$ juga
!!{axiomatizable}. Untuk menguji apakah $!B$ termasuk aksioma~$\Th{PA}$,
yakni untuk menghitung fungsi $\Char{X}$ dengan $\Char{X}(!B) = 1$ jika $!B$
adalah aksioma~$\Th{PA}$ dan $= 0$ jika tidak, kita dapat melakukan hal
berikut. Pertama, periksa apakah $!B$ merupakan salah satu aksioma~$\Th{Q}$.
Jika ya, jawabannya ``ya'' dan nilai $\Char{X}(!B) = 1$. Jika tidak, uji apakah
kalimat tersebut merupakan contoh skema induksi. Pengujian ini dapat dilakukan secara
sistematis; dalam hal ini, barangkali cara termudah untuk melihatnya ialah
sebagai berikut. Setiap contoh skema induksi dimulai dengan sejumlah kuantor
universal, lalu suatu sub-!!{formula} yang berbentuk implikasi. Konsekuen
implikasi itu ialah $\lforall[x][!A(x, y_1, \dots, y_n)]$, dengan $x$ dan $y_1$,
\dots, $y_n$ sebagai semua variabel bebas~$!A$, sedangkan kuantor awal~$!B$
mengikat variabel $y_1$, \dots, $y_n$. Setelah mengekstrak~$!A$ ini dan
memeriksa bahwa variabel bebasnya cocok dengan variabel yang diikat oleh
kuantor universal di bagian depan dan oleh~$\lforall[x]$, kita memeriksa
bahwa anteseden implikasinya cocok dengan
\[
!A(\Obj 0, y_1, \dots, y_n) \land \lforall[x][(!A(x, y_1, \dots, y_n)
\lif !A(x', y_1, \dots, y_n))]
\]
Sekali lagi, jika cocok, $!B$ merupakan contoh skema induksi; jika tidak,
$!B$ bukan contoh skema induksi.
\end{ex}

Dalam menjawab pertanyaan ini---dan pertanyaan yang lebih umum tentang teori
mana yang lengkap atau dapat diputuskan---kita juga akan terbantu oleh
definisi berikut. Ingat bahwa suatu himpunan $X$ !!{enumerable} jika dan hanya
jika kosong atau terdapat !!a{surjective} fungsi~$f \colon \Nat \to X$.
Fungsi semacam itu disebut enumerasi~$X$.

\begin{defn}
Suatu himpunan $X$ disebut \emph{!!{computably enumerable}} (disingkat
!!{c.e.}) jika dan hanya jika kosong atau memiliki enumerasi yang dapat
dihitung.
\end{defn}

Selain !!{axiomatizability}, syarat lain bagi teori yang terkena
teorema-teorema ketaklengkapan ialah bahwa teori itu cukup kuat untuk
membuktikan fakta dasar tentang fungsi yang dapat dihitung dan relasi
!!{decidable}. Dengan ``fakta dasar,'' kita maksudkan !!{sentence}s yang
menyatakan nilai fungsi yang dapat dihitung untuk masing-masing argumennya.
Dengan ``cukup kuat,'' kita maksudkan bahwa teori yang bersangkutan memuat
kalimat-kalimat itu sebagai teorema. Sebagai contoh, perhatikan fungsi yang
dapat dihitung paling lazim: penjumlahan. Nilai $+$ untuk argumen $2$ dan $3$
ialah~$5$, yakni $2+3 = 5$. Kalimat dalam bahasa aritmetika yang menyatakan
bahwa nilai $+$ untuk argumen $2$ dan $3$ ialah~$5$ adalah
$(\num{2} + \num{3}) = \num{5}$. Sebagai contoh, $\Th{Q}$ membuktikan kalimat
ini. Secara lebih umum, untuk setiap fungsi yang dapat dihitung
$f(x_1, x_2)$, kita menginginkan !!a{formula}~$!A_f(x_1, x_2, y)$ dalam
$\Lang{L_A}$ sedemikian sehingga $\Th{Q} \Proves
!A_f(\num{n_1}, \num{n_2}, \num{m})$ setiap kali $f(n_1, n_2) = m$. Dengan cara
ini, $\Th{Q}$ membuktikan bahwa nilai~$f$ untuk argumen $n_1$, $n_2$ ialah~$m$.
Kita bahkan memerlukan sedikit lebih banyak: teori itu harus membuktikan bahwa
tidak ada bilangan lain yang merupakan nilai~$f$ untuk argumen $n_1$,~$n_2$.
Hal yang sama berlaku bagi relasi !!{decidable}. Kedua definisi berikut
membuatnya tepat.

\begin{defn}
Suatu !!{formula}~$!A(x_1, \dots, x_k, y)$ \emph{!!{represents}} fungsi
$f\colon \Nat^k \to \Nat$ dalam~$\Gamma$ jika dan hanya jika, setiap kali
$f(n_1, \dots, n_k) = m$, berlaku
\begin{enumerate}
\item $\Gamma \Proves !A(\num{n_1}, \dots, \num{n_k}, \num{m})$, dan
\item $\Gamma \Proves \lforall[y](!A(\num{n_1}, \dots, \num{n_k},
y) \lif y = \num{m})$.
\end{enumerate}
\end{defn}

\begin{defn}
Suatu !!{formula}~$!A(x_1, \dots, x_k)$ \emph{!!{represents}} relasi
$R \subseteq \Nat^k$ dalam~$\Gamma$ jika dan hanya jika
\begin{enumerate}
\item setiap kali $R(n_1, \dots, n_k)$, $\Gamma \Proves
  !A(\num{n_1}, \dots, \num{n_k})$, dan
\item setiap kali $R(n_1, \dots, n_k)$ tidak berlaku, $\Gamma \Proves \lnot
  !A(\num{n_1}, \dots, \num{n_k})$.
\end{enumerate}
\end{defn}

Suatu teori ``cukup kuat'' agar teorema ketaklengkapan berlaku jika teori itu
!!{represents} semua fungsi yang dapat dihitung dan semua relasi
!!{decidable}. $\Th{Q}$ beserta perluasannya memenuhi syarat ini, tetapi
memerlukan waktu untuk menetapkannya---ini merupakan fakta tak sepele tentang
hal-hal yang dapat dibuktikan $\Th{Q}$, dan sulit ditunjukkan karena $\Th{Q}$
hanya memiliki sedikit aksioma sebagai dasar pembuktian. Namun, $\Th{Q}$
adalah teori yang sangat lemah. Jadi, meskipun sulit membuktikan bahwa
$\Th{Q}$ merepresentasikan semua fungsi yang dapat dihitung, kebanyakan teori
yang menarik lebih kuat daripada~$\Th{Q}$, yakni membuktikan lebih banyak
daripada $\Th{Q}$. Apa pun yang dibuktikan $\Th{Q}$ juga dibuktikan oleh
setiap teori yang lebih kuat; karena $\Th{Q}$ merepresentasikan semua fungsi
yang dapat dihitung, setiap teori yang lebih kuat juga demikian. Artinya,
banyak teori menarik memenuhi syarat teorema ketaklengkapan ini. Kerja keras
kita akan terbayar karena menunjukkan bahwa teorema ketaklengkapan berlaku
bagi berbagai macam teori. Tentu saja, setiap teori yang bermaksud
memformalkan ``seluruh matematika'' harus membuktikan semua yang dibuktikan
$\Th{Q}$, sebab setidaknya teori tersebut harus mampu menangkap hasil
komputasi elementer. Jadi, setiap calon teori ``seluruh matematika'' termasuk
teori yang terkena teorema ketaklengkapan.

\end{document}
